\documentclass{apcalc}
\title{Unit 3 Practice Set: Differentiation: Composite, Implicit, and Inverse Functions}
\author{Paul Wang}
\date{November 2026}
\volume{Practice Set 3 of 10 (topics 3.1–3.6)}
\shorttitle{Unit 3 --- Practice Set}
\begin{document}
\maketitle
\begin{abstract}
Twelve exam-style multiple-choice questions on Unit 3, grouped by the sections of the Unit 3 Study Guide, and one nine-point free-response question on an implicitly defined curve. All items are no-calculator. Write each answer in the blank beside the number; the key on the last pages gives an answers-only strip, a reason and rule for every item, and the free-response scoring guide.
\end{abstract}
\scoringline{12 MC \quad\blank\ / 9 FR}{Exam pace: about 2 minutes per multiple-choice item (24 min) and 15 minutes for the free response.}
\section*{1\quad The chain rule}
\begin{setrule}R3.1–R3.2. $\frac{d}{dx}f(g(x))=f'(g(x))\,g'(x)$. With a table, find $g(a)$ first, then $f'$ at that output.\end{setrule}
\begin{q}{1}
If $f(x)=(x^2+1)^5$, then $f'(1)=$
\chfour{$10$}{$32$}{$80$}{$160$}
\end{q}
\begin{q}{2}
$\dfrac{d}{dx}\ln(\cos x)=$
\chfour{$-\tan x$}{$\dfrac{1}{\cos x}$}{$-\dfrac{1}{\sin x}$}{$\tan x$}
\end{q}
\begin{q}{3}
The table gives values of $f$, $g$ and their derivatives. If $h(x)=f(g(x))$, then $h'(1)=$
\begin{center}\small\renewcommand{\arraystretch}{1.35}\begin{tabular}{|c|c|c|c|c|}
\hline
$x$ & $f(x)$ & $f'(x)$ & $g(x)$ & $g'(x)$ \\ \hline
1 & 4 & $-1$ & 3 & 2 \\ \hline
3 & 2 & 6 & 1 & $-2$ \\ \hline
\end{tabular}\end{center}
\chfour{$-2$}{$6$}{$12$}{$18$}
\end{q}
\begin{q}{4}
With the same table, if $k(x)=[g(x)]^3$, then $k'(3)=$
\begin{center}\small\renewcommand{\arraystretch}{1.35}\begin{tabular}{|c|c|c|c|c|}
\hline
$x$ & $f(x)$ & $f'(x)$ & $g(x)$ & $g'(x)$ \\ \hline
1 & 4 & $-1$ & 3 & 2 \\ \hline
3 & 2 & 6 & 1 & $-2$ \\ \hline
\end{tabular}\end{center}
\chfour{$-6$}{$-2$}{$3$}{$6$}
\end{q}
\section*{2\quad Implicit differentiation}
\begin{setrule}R3.3–R3.4. Every $y$-term gets a factor $\frac{dy}{dx}$. Horizontal tangent: numerator 0; vertical: denominator 0; the point must be on the curve.\end{setrule}
\begin{q}{5}
For the curve $x^2y+y^3=10$, the value of $\dfrac{dy}{dx}$ at the point $(1,2)$ is
\chfour{$-16$}{$-\dfrac13$}{$-\dfrac{4}{13}$}{$\dfrac{4}{13}$}
\end{q}
\begin{q}{6}
At which points does the curve $x^2-xy+y^2=3$ have a horizontal tangent line?
\chtwo{$(2,1)$ and $(-2,-1)$}{$(1,2)$ and $(-1,-2)$}{$(\sqrt3,0)$ and $(-\sqrt3,0)$}{$(0,\sqrt3)$ and $(0,-\sqrt3)$}
\end{q}
\section*{3\quad Derivatives of inverse functions}
\begin{setrule}R3.5–R3.6. $(f^{-1})'(b)=\frac{1}{f'(a)}$ where $f(a)=b$. $\arcsin'=\frac{1}{\sqrt{1-x^2}}$, $\arctan'=\frac{1}{1+x^2}$.\end{setrule}
\begin{q}{7}
Let $f(x)=x^3+2x-1$ and let $g$ be the inverse of $f$. What is $g'(2)$?
\chfour{$\dfrac{1}{14}$}{$\dfrac15$}{$5$}{$14$}
\end{q}
\begin{q}{8}
A differentiable, invertible $f$ has $f(3)=7$, $f'(3)=4$, $f(7)=10$ and $f'(7)=-\frac12$. If $g=f^{-1}$, then $g'(7)=$
\chfour{$-2$}{$\dfrac{1}{10}$}{$\dfrac14$}{$4$}
\end{q}
\begin{q}{9}
If $y=\arctan(3x)$, then $\dfrac{dy}{dx}$ at $x=\frac13$ is
\chfour{$\dfrac12$}{$\dfrac32$}{$\dfrac94$}{$3$}
\end{q}
\begin{q}{10}
$\dfrac{d}{dx}\arcsin(x^2)=$
\chfour{$\dfrac{1}{\sqrt{1-x^4}}$}{$\dfrac{2x}{\sqrt{1-x^2}}$}{$-\dfrac{2x}{\sqrt{1-x^4}}$}{$\dfrac{2x}{\sqrt{1-x^4}}$}
\end{q}
\section*{4\quad Selecting procedures and higher-order derivatives}
\begin{setrule}R3.7–R3.8. Identify the last operation. For an implicit $\frac{d^2y}{dx^2}$, differentiate $\frac{dy}{dx}$ and substitute it back.\end{setrule}
\begin{q}{11}
For the curve $y^2=x^3$ with $y>0$, $\dfrac{d^2y}{dx^2}$ at the point $(1,1)$ is
\chfour{$\dfrac34$}{$\dfrac32$}{$3$}{$\dfrac{21}{4}$}
\end{q}
\begin{q}{12}
If $f(x)=\sin(2x)$, then $f^{(4)}(x)$, the fourth derivative, is
\chtwo{$-16\sin(2x)$}{$8\cos(2x)$}{$16\cos(2x)$}{$16\sin(2x)$}
\end{q}
\clearpage\section*{Free response}
Show all work. Justify every conclusion with calculus, and give decimal answers correct to three places.
\begin{frq}{Question 1 --- An implicitly defined curve (9 points, no calculator)}
Consider the curve given by $y^2+xy-x^2=5$.
\partlabel{a}Show that $\dfrac{dy}{dx}=\dfrac{2x-y}{2y+x}$.\workspace{1.6in}
\partlabel{b}The point $(1,-3)$ is on the curve. Write an equation for the line tangent to the curve at this point.\workspace{1.6in}
\partlabel{c}Show that the curve has no vertical tangent lines.\workspace{1.6in}
\partlabel{d}Find the value of $\dfrac{d^2y}{dx^2}$ at the point $(1,-3)$.\workspace{1.6in}
\end{frq}
\clearpage\section*{Answer key}
\keystrip{\keyitem{1}{D}\keyitem{2}{A}\keyitem{3}{C}\keyitem{4}{A}\keyitem{5}{C}\keyitem{6}{B}\keyitem{7}{B}\keyitem{8}{C}\keyitem{9}{B}\keyitem{10}{D}\keyitem{11}{A}\keyitem{12}{D}}
\vspace{4pt}
\begin{keytable}
\keysection{1. The chain rule}
1 & D & $5(x^2+1)^4\cdot2x=5\cdot16\cdot2=160$. & R3.1\\
2 & A & $\frac{-\sin x}{\cos x}=-\tan x$. & R3.1\\
3 & C & $f'(g(1))g'(1)=f'(3)\cdot2=12$. & R3.2\\
4 & A & $3g(3)^2g'(3)=3\cdot1\cdot(-2)=-6$. & R3.2\\
\keysection{2. Implicit differentiation}
5 & C & $2xy+x^2y'+3y^2y'=0\Rightarrow y'=\frac{-2xy}{x^2+3y^2}=-\frac{4}{13}$. & R3.3\\
6 & B & $y'=\frac{y-2x}{2y-x}$; $y=2x$ on the curve gives $3x^2=3$. & R3.4\\
\keysection{3. Derivatives of inverse functions}
7 & B & $f(1)=2$, so $g'(2)=\frac{1}{f'(1)}=\frac15$. & R3.5\\
8 & C & $g(7)=3$, so $g'(7)=\frac{1}{f'(3)}=\frac14$. & R3.5\\
9 & B & $\frac{3}{1+(3x)^2}=\frac{3}{2}$. & R3.6\\
10 & D & $\frac{u'}{\sqrt{1-u^2}}$ with $u=x^2$. & R3.6\\
\keysection{4. Selecting procedures and higher-order derivatives}
11 & A & $y'=\frac{3x^2}{2y}=\frac32$; $y''=\frac{6x(2y)-3x^2(2y')}{4y^2}=\frac{12-9}{4}$. & R3.8\\
12 & D & Each derivative multiplies by 2 and cycles $\sin\to\cos\to-\sin\to-\cos\to\sin$. & R3.7\\
\end{keytable}
\section*{Free-response scoring guide}
\subsection*{Question 1: An implicitly defined curve (9 points)}
\textbf{(a)} $2y\frac{dy}{dx}+y+x\frac{dy}{dx}-2x=0$, so $\frac{dy}{dx}(2y+x)=2x-y$ and $\frac{dy}{dx}=\frac{2x-y}{2y+x}$.
\begin{itemize}[leftmargin=1.6em,itemsep=0pt]
\item 1 point: implicit differentiation of all three terms (with the product rule on $xy$)
\item 1 point: algebra leading to the given expression
\end{itemize}
\textbf{(b)} At $(1,-3)$: $\frac{dy}{dx}=\frac{2+3}{-6+1}=-1$. Tangent: $y=-3-(x-1)$.
\begin{itemize}[leftmargin=1.6em,itemsep=0pt]
\item 1 point: slope $-1$
\item 1 point: tangent line $y=-3-(x-1)$
\end{itemize}
\textbf{(c)} A vertical tangent needs $2y+x=0$, so $x=-2y$. Then $y^2-2y^2-4y^2=-5y^2=5$, so $y^2=-1$, which has no real solution. No point of the curve has $2y+x=0$, so there is no vertical tangent.
\begin{itemize}[leftmargin=1.6em,itemsep=0pt]
\item 1 point: sets $2y+x=0$
\item 1 point: substitutes $x=-2y$ into the curve: $-5y^2=5$
\item 1 point: concludes: no real solution, so no vertical tangent
\end{itemize}
\textbf{(d)} $\frac{d^2y}{dx^2}=\frac{(2-y')(2y+x)-(2x-y)(2y'+1)}{(2y+x)^2}$. At $(1,-3)$ with $y'=-1$: $\frac{(3)(-5)-(5)(-1)}{25}=\frac{-10}{25}=-\frac25$.
\begin{itemize}[leftmargin=1.6em,itemsep=0pt]
\item 1 point: quotient rule applied to $\frac{dy}{dx}$ with $\frac{dy}{dx}$ retained
\item 1 point: value $-\frac25$
\end{itemize}
\end{document}